\section{第10章\quad 重积分}
\subsection{二重积分的概念与性质}

\begin{exercise}{284}{基础·积分方程}
设 $D=[0,1]^2$，连续函数满足 $f(x,y)=xy+\frac12\iint_D f(u,v)\,\dd u\dd v$。令积分为 $C$，求 $C$ 及 $f(x,y)$。
\end{exercise}
\begin{exercise}{285}{基础·二重和极限}
将 $\lim_{n\to\infty}n^{-2}\sum_{i=1}^n\sum_{j=1}^n(i/n+j/n)$ 化为单位正方形上的二重积分并计算。
\end{exercise}
\begin{exercise}{286}{基础·黎曼和}
计算 $\lim_{n\to\infty}\frac{\pi}{2n^4}\sum_{i=1}^n\sum_{j=1}^n i^2\sin\frac{j\pi}{2n}$。提示：写成 $x^2\sin y$ 的二重积分。
\end{exercise}
\begin{exercise}{287}{基础·积分比较}
设 $D=\{x,y\geq0:x+y\leq1\}$。比较 $\iint_D(x+y)^2\,\dd x\dd y$ 与 $\iint_D(x+y)\,\dd x\dd y$ 的大小，说明依据。
\end{exercise}
\begin{exercise}{288}{基础·函数大小与积分}
设 $D=[0,1]^2$，$u=(x+y)/2$。按从小到大排列 $\iint_Du\,\dd\sigma,\iint_D\sqrt u\,\dd\sigma,\iint_D\sqrt[3]u\,\dd\sigma$，无需计算积分值。
\end{exercise}
\begin{exercise}{289}{基础·变量分离}
设 $D=[0,\pi]^2$，求 $\iint_D\sin^2x\sin^2y\,\dd x\dd y$ 及被积函数在 $D$ 上的平均值。
\end{exercise}
\begin{exercise}{290}{基础·积分估值}
设 $D=\{x^2+y^2\leq1\}$。利用 $1\leq1+x^2+y^2\leq2$ 给出 $\iint_D(1+x^2+y^2)\,\dd\sigma$ 的上下界。
\end{exercise}
\begin{exercise}{291}{基础·分片常值积分}
设 $D=[0,3]^2$，$f$ 在 $0\leq y\leq1$ 上为 $1$，在 $1<y\leq2$ 上为 $2$，在 $2<y\leq3$ 上为 $3$。求 $\iint_D f\,\dd\sigma$。
\end{exercise}
\begin{exercise}{292}{基础·积分平均值}
设 $D=\{0\leq y\leq x\leq1\}$。求 $xy$ 在 $D$ 上的平均值。
\end{exercise}
\begin{exercise}{293}{基础·小圆域平均值}
设 $D_r=\{x^2+y^2\leq r^2\}$。利用积分中值定理求 $\lim_{r\to0^+}\frac1{\pi r^2}\iint_{D_r}e^{x+y}\,\dd x\dd y$。
\end{exercise}
\begin{exercise}{294}{基础·连续函数积分极限}
设 $f$ 在原点附近连续且 $f(0,0)=2$。求 $\lim_{r\to0^+}r^{-2}\iint_{x^2+y^2\leq r^2}f(x,y)\,\dd x\dd y$。
\end{exercise}
\begin{exercise}{295}{基础·奇偶对称性}
计算 $\iint_{[-1,1]^2}(x^3+xy+y^2)\,\dd x\dd y$，指出哪些项由对称性消去。
\end{exercise}
\begin{exercise}{296}{常规·变量互换}
设 $D=\{x^2+y^2\leq1,x,y\geq0\}$。利用 $x,y$ 互换计算 $\iint_D\frac{e^x}{e^x+e^y}\,\dd x\dd y$。
\end{exercise}
\begin{exercise}{297}{基础·三角形积分}
设 $D=\{-1\leq x\leq y\leq1\}$。计算 $\iint_D(x+y)\,\dd x\dd y$。
\end{exercise}
\begin{exercise}{298}{基础·圆域对称性}
设 $D=\{x^2+y^2\leq1\}$。利用对称性与极坐标计算 $\iint_D(x^2+2y)\,\dd x\dd y$。
\end{exercise}
\begin{exercise}{299}{基础·非负函数积分}
设 $D=[0,1]^2$，$f$ 在 $D$ 上连续且非负。证明若 $\iint_D f\,\dd\sigma=0$，则 $f\equiv0$。提示：若某点函数值为正，利用连续性在其附近给出正的下界。
\end{exercise}
\begin{exercise}{300}{常规·积分方程小综合}
设 $D=[0,1]^2$，连续函数满足 $f(x,y)=x+y+\frac12\iint_D f(u,v)\,\dd u\dd v$。求 $f$ 并计算它在 $D$ 上的平均值。
\end{exercise}

\subsection{二重积分的计算方法}
\subsubsection{积分区域与交换积分次序}
\begin{exercise}{301}{基础·交换积分次序}
将 $\int_0^1\dd x\int_x^1 f(x,y)\,\dd y$ 交换积分次序，其中 $f$ 连续；取 $f=y$ 计算积分。
\end{exercise}
\begin{exercise}{302}{基础·半圆域换序}
将 $\int_0^1\dd y\int_{-\sqrt{1-y^2}}^{\sqrt{1-y^2}}f(x,y)\,\dd x$ 交换次序，其中 $f$ 在闭半圆上连续。
\end{exercise}
\begin{exercise}{303}{基础·分块区域合并}
设 $f$ 连续。把 $\int_0^1\dd x\int_0^x f\,\dd y+\int_1^2\dd x\int_0^{2-x}f\,\dd y$ 改写为一个先对 $x$ 积分的累次积分。
\end{exercise}
\begin{exercise}{304}{基础·三角形换序}
设 $f$ 连续。将 $\int_0^1\dd y\int_0^{2y}f\,\dd x+\int_1^2\dd y\int_0^{4-2y}f\,\dd x$ 交换次序。
\end{exercise}
\begin{exercise}{305}{基础·曲线区域面积}
设 $D=\{0\leq x\leq\pi,0\leq y\leq\sin x\}$。写出先对 $x$ 积分的面积积分并求 $D$ 的面积。
\end{exercise}
\begin{exercise}{306}{常规·抛物线域换序}
设 $f$ 连续。将 $\int_1^2\dd y\int_y^{y^2}f(x,y)\,\dd x$ 交换次序，写明需要分段的点。
\end{exercise}
\begin{exercise}{307}{基础·区域重新组合}
设 $f$ 连续。将 $\int_0^1\dd x\int_0^x f\,\dd y+\int_1^2\dd x\int_0^1 f\,\dd y$ 改为先对 $x$ 积分的累次积分。
\end{exercise}
\begin{exercise}{308}{基础·变限二重积分}
设 $f$ 连续，$F(t)=\int_0^t\dd y\int_y^t f(x)\,\dd x$，$t\geq0$。先换序，再求 $F'(t)$。
\end{exercise}
\subsubsection{直角坐标计算与分区}
\begin{exercise}{309}{基础·矩形域计算}
计算 $\iint_{[0,1]^2}e^{x-y}\,\dd x\dd y$。
\end{exercise}
\begin{exercise}{310}{基础·区域可加性}
设 $D=[0,1]^2$，$D_-=D\cap\{x+y\leq1\}$。求 $x+y$ 在 $D_-$ 及 $D$ 上的积分。
\end{exercise}
\begin{exercise}{311}{基础·指数积分}
计算 $\int_0^1\dd x\int_0^1xe^{xy}\,\dd y$。
\end{exercise}
\begin{exercise}{312}{基础·三角形域计算}
设 $D=\{1\leq y\leq x\leq2\}$。求 $\iint_Dxy\,\dd x\dd y$。
\end{exercise}
\begin{exercise}{313}{基础·曲边薄片质量}
设 $D$ 由 $y=\sqrt{x}$ 与 $y=x^2$ 围成，面密度为 $\rho=1$。求其质量并写出交点。
\end{exercise}
\begin{exercise}{314}{常规·三角形质心}
设 $D=\{x,y\geq0,x+y\leq1\}$，面密度为 $\rho=x+y$。求质量及质心坐标。
\end{exercise}
\begin{exercise}{315}{基础·三角函数积分}
设 $D=\{0\leq y\leq x\leq\pi\}$。求 $\iint_D\cos(x+y)\,\dd x\dd y$。
\end{exercise}
\begin{exercise}{316}{基础·双曲线边界}
设 $D=\{1\leq x\leq2,1/x\leq y\leq x\}$。求 $\iint_D x\,\dd x\dd y$。
\end{exercise}
\begin{exercise}{317}{基础·非均匀薄片质量}
设 $D=\{0\leq x\leq1,x\leq y\leq2x\}$，密度为 $\rho=x+y$。求质量。
\end{exercise}
\begin{exercise}{318}{基础·梯形域积分}
设 $D=\{0\leq x\leq1,0\leq y\leq1+x\}$。求 $\iint_D(1+x)\,\dd x\dd y$。
\end{exercise}
\begin{exercise}{319}{基础·多项式积分}
设 $D=\{0\leq y\leq2,y/2\leq x\leq y\}$。计算 $\iint_D(x^2+y^2)\,\dd x\dd y$。
\end{exercise}
\begin{exercise}{320}{基础·平行四边形换元}
设 $D=\{1\leq y\leq3,y-1\leq x\leq y\}$。用 $u=y-x,v=y$ 计算 $\iint_D(y-x)\,\dd x\dd y$，写出雅可比绝对值。
\end{exercise}
\begin{exercise}{321}{基础·三角形形心}
设均匀薄片的区域为顶点 $(0,0),(1,2),(2,1)$ 所成的三角形。用二重积分求形心。
\end{exercise}
\begin{exercise}{322}{基础·积分次序选择}
设 $D=\{0\leq y\leq1,0\leq x\leq y\}$。选择适当次序计算 $\iint_D\cos(y^2)\,\dd x\dd y$。
\end{exercise}
\begin{exercise}{323}{常规·可去点与积分}
令 $f(x)=\sin x/x$（$x\ne0$），$f(0)=1$。设 $D=\{0\leq x\leq1,0\leq y\leq x\}$，计算 $\iint_Df(x)\,\dd x\dd y$。
\end{exercise}
\begin{exercise}{324}{基础·换序计算}
计算 $\int_0^1\dd x\int_x^1 e^{-y^2}\,\dd y$。
\end{exercise}
\begin{exercise}{325}{基础·含高斯因子的积分}
设 $D=\{0\leq x\leq y\leq1\}$。计算 $\iint_Dx^2e^{-y^2}\,\dd x\dd y$。提示：先对 $x$ 积分，再令 $u=y^2$。
\end{exercise}
\begin{exercise}{326}{基础·区域宽度}
设 $D=\{0\leq x\leq1,0\leq y\leq x\}$。求 $\iint_De^{x^2}\,\dd x\dd y$。
\end{exercise}
\begin{exercise}{327}{基础·换序与分离}
将 $\int_0^1\dd x\int_x^1(y+xye^{y^2})\,\dd y$ 交换积分次序并计算。
\end{exercise}
\begin{exercise}{328}{基础·换序与代换}
计算 $\int_0^1\dd x\int_x^1\frac{x}{\sqrt{1+y^3}}\,\dd y$。提示：先换序。
\end{exercise}
\begin{exercise}{329}{基础·对数边界积分}
设 $D=\{1\leq x\leq e,0\leq y\leq\ln x\}$。计算 $\iint_Dx^{-1}\,\dd x\dd y$。
\end{exercise}
\begin{exercise}{330}{常规·换序与分部积分}
设 $f\in C^1[0,1]$，$f(0)=0,f(1)=1,\int_0^1 f(y)\,\dd y=1/3$。先换序，再计算 $\int_0^1\dd x\int_0^x f'(y)\,\dd y$。
\end{exercise}
\begin{exercise}{331}{常规·支集分段}
设 $f(x)=1$（$0\leq x\leq1$），其余处为 $0$。求 $H(y)=\int_{-\infty}^{+\infty}f(x)f(y-x)\,\dd x$ 的分段表达式。
\end{exercise}
\begin{exercise}{332}{基础·支集与区域}
设 $f(x,y)=xy$ 在 $[0,1]^2$ 上定义，其余处为零。求 $\iint_{x+y\leq1}f(x,y)\,\dd x\dd y$。
\end{exercise}
\begin{exercise}{333}{基础·绝对值分区}
设 $D=[0,1]^2$。沿 $y=x$ 分区计算 $\iint_D|y-x|\,\dd x\dd y$。
\end{exercise}
\begin{exercise}{334}{基础·最大值函数}
设 $D=[0,1]^2$。计算 $\iint_D\max\{x,y\}\,\dd x\dd y$，写出分区界线。
\end{exercise}
\begin{exercise}{335}{常规·对称分区积分}
设 $D=[0,1]^2$。利用 $x,y$ 互换，计算 $\iint_De^{\max\{x^2,y^2\}}\,\dd x\dd y$。
\end{exercise}
\begin{exercise}{336}{常规·根式边界换元}
设 $D=\{x,y\geq0,\sqrt x+\sqrt y\leq1\}$。用 $u=\sqrt x,v=\sqrt y$ 求 $D$ 的面积，写出换元后的积分区域。
\end{exercise}
\subsubsection{极坐标与圆形边界}
\begin{exercise}{337}{基础·扇环域}
设 $D=\{(r\cos\theta,r\sin\theta):1\leq r\leq2,\pi/6\leq\theta\leq\pi/3\}$。求 $\iint_D\arctan(y/x)\,\dd x\dd y$。
\end{exercise}
\begin{exercise}{338}{基础·圆域根式积分}
设 $D=\{x^2+y^2\leq1,x,y\geq0\}$。用极坐标计算 $\iint_D\sqrt{1-x^2-y^2}\,\dd x\dd y$。
\end{exercise}
\begin{exercise}{339}{基础·半球体积}
设 $a>0$。计算 $\iint_{x^2+y^2\leq a^2}\sqrt{a^2-x^2-y^2}\,\dd x\dd y$，说明它代表的体积。
\end{exercise}
\begin{exercise}{340}{基础·圆盘平均距离}
设 $D=\{x^2+y^2\leq1,x,y\geq0\}$。求 $\iint_D\sqrt{x^2+y^2}\,\dd x\dd y$ 及距离的平均值。
\end{exercise}
\begin{exercise}{341}{基础·平移圆域}
设 $D=\{x^2+(y-1)^2\leq1\}$。作平移换元计算 $\iint_D y\,\dd x\dd y$。
\end{exercise}
\begin{exercise}{342}{基础·环域对数积分}
设 $D=\{1\leq x^2+y^2\leq4\}$。用极坐标计算 $\iint_D\ln(x^2+y^2)\,\dd x\dd y$。
\end{exercise}
\begin{exercise}{343}{常规·对称性与极坐标}
设 $D=\{x^2+y^2\leq1,x,y\geq0\}$，原点处将被积函数定义为零。利用变量互换计算 $\iint_D\frac{x\cos(x^2+y^2)}{x+y}\,\dd x\dd y$。
\end{exercise}
\begin{exercise}{344}{基础·径向反常积分}
用圆心截断 $\varepsilon\leq\sqrt{x^2+y^2}\leq1$ 定义并计算 $\iint_{x^2+y^2\leq1}(x^2+y^2)^{-1/2}\,\dd x\dd y$。
\end{exercise}
\begin{exercise}{345}{基础·扇形径向积分}
设 $D=\{x^2+y^2\leq4,0\leq y\leq x\}$。用极坐标计算 $\iint_D\sqrt{x^2+y^2}\,\dd x\dd y$。
\end{exercise}
\begin{exercise}{346}{常规·三角形反常积分}
计算 $\int_0^1\dd x\int_x^{\sqrt3x}\frac1{\sqrt{x^2+y^2}}\,\dd y$，原点按圆心截断极限理解。提示：极坐标中 $r\cos\theta\leq1$。
\end{exercise}
\begin{exercise}{347}{基础·偏心圆域}
设 $D=\{x^2+y^2\leq2x,y\geq0\}$。写出极坐标范围并求 $D$ 的面积。
\end{exercise}
\begin{exercise}{348}{基础·半圆域对称性}
设 $D=\{x^2+y^2\leq1,y\geq0\}$。求 $\iint_D(x+y)\,\dd x\dd y$。
\end{exercise}
\begin{exercise}{349}{基础·偏心圆的矩}
设 $D=\{(x-1)^2+y^2\leq1\}$。利用平移及对称性求 $\iint_D(x+y^2)\,\dd x\dd y$。
\end{exercise}
\begin{exercise}{350}{常规·圆域积分方程}
设 $D=\{x^2+y^2\leq1\}$，连续函数满足 $f(x,y)=x^2+y^2+\frac1{2\pi}\iint_Df\,\dd\sigma$。求 $f$。
\end{exercise}
\begin{exercise}{351}{基础·径向分区}
设 $D=\{x^2+y^2\leq1\}$。计算 $\iint_D\min\{x^2+y^2,1/4\}\,\dd x\dd y$，写出分界圆。
\end{exercise}
\begin{exercise}{352}{基础·绝对值与对称性}
设 $D=\{x^2+y^2\leq1\}$。计算 $\iint_D|xy|\,\dd x\dd y$。
\end{exercise}
\subsubsection{一般换元与无界积分}
\begin{exercise}{353}{基础·线性换元}
设 $D=\{0\leq x+y\leq1,0\leq x-y\leq1\}$。令 $u=x+y,v=x-y$，计算 $\iint_D(x+y)\,\dd x\dd y$，写出雅可比绝对值。
\end{exercise}
\begin{exercise}{354}{基础·椭圆换元}
设 $D=\{x^2/4+y^2\leq1\}$。用 $x=2r\cos\theta,y=r\sin\theta$ 计算 $\iint_D(x^2/4+y^2)\,\dd x\dd y$。
\end{exercise}
\begin{exercise}{355}{常规·和与比值换元}
设 $D=\{x,y\geq0,x+y\leq1\}$，原点处赋值为 $0$。用 $u=x+y,v=x/(x+y)$ 计算 $\iint_De^{x/(x+y)}\,\dd x\dd y$；写出逆变换及雅可比绝对值。
\end{exercise}
\begin{exercise}{356}{常规·乘积换元}
设 $D=\{x,y>0,1\leq xy\leq2,1\leq y/x\leq4\}$。用 $u=xy,v=y/x$ 求 $D$ 的面积。提示：$x=\sqrt{u/v},y=\sqrt{uv}$。
\end{exercise}
\begin{exercise}{357}{基础·菱形域换元}
设 $D=\{|x|+|y|\leq1\}$。用 $u=x+y,v=x-y$ 计算 $\iint_D(x+y)^2\,\dd x\dd y$。
\end{exercise}
\begin{exercise}{358}{基础·无界区域换序}
设 $D=\{x\geq0,x^2\leq y\leq4x^2\}$。先以 $y\leq R$ 截断并交换次序，再求 $\iint_Dxe^{-y}\,\dd x\dd y$。
\end{exercise}
\begin{exercise}{359}{基础·高斯型积分}
用半径 $R$ 的圆盘截断，计算第一象限上的 $\iint e^{-(x^2+y^2)}\,\dd x\dd y$。
\end{exercise}
\begin{exercise}{360}{基础·无界径向积分}
用第一象限半径 $R$ 的圆盘递增截断，计算 $\iint_{x,y\geq0}(1+x^2+y^2)^{-2}\,\dd x\dd y$。
\end{exercise}
\begin{exercise}{361}{基础·换元区域变式}
设 $D=\{0\leq x+y\leq2,0\leq x-y\leq1\}$。写出 $u=x+y,v=x-y$ 的逆变换，并用换元计算 $D$ 的面积。
\end{exercise}
\begin{exercise}{362}{常规·反常积分小综合}
设 $D_\varepsilon=\{\varepsilon^2\leq x^2+y^2\leq1\}$。计算 $\iint_{D_\varepsilon}(x^2+y^2)^{-1}\,\dd x\dd y$，判断 $\varepsilon\to0^+$ 时积分是否有有限极限，并与幂次 $-1/2$ 的情形比较。
\end{exercise}
\subsection{三重积分}
\begin{exercise}{363}{基础·直角坐标三重积分}
设 $\Omega=\{0\leq y\leq x\leq1,0\leq z\leq1\}$。写出累次积分并计算 $\iiint_\Omega xy\,\dd x\dd y\dd z$。
\end{exercise}
\begin{exercise}{364}{基础·积分次序选择}
设 $\Omega=\{0\leq x\leq\pi/2,0\leq y\leq x,0\leq z\leq1\}$。计算 $\iiint_\Omega\sin x\,\dd V$。
\end{exercise}
\begin{exercise}{365}{基础·椭球截面}
设 $\Omega=\{x^2+y^2/4+z^2/9\leq1,z\geq0\}$。用水平截面计算 $\iiint_\Omega z\,\dd V$。
\end{exercise}
\begin{exercise}{366}{基础·抛物面围体积}
设 $\Omega=\{x^2+y^2\leq z\leq1\}$。写出柱面坐标范围并求体积。
\end{exercise}
\begin{exercise}{367}{基础·水平截层积分}
设 $\Omega=\{1\leq z\leq2,x^2+y^2\leq z\}$。求 $\iiint_\Omega(x^2+y^2)\,\dd V$。
\end{exercise}
\begin{exercise}{368}{基础·旋转曲面围体}
将 $x=0$ 平面内的抛物线 $y^2=2z$ 绕 $z$ 轴旋转，与 $z=1$ 围成 $\Omega$。写出曲面方程并求体积。
\end{exercise}
\begin{exercise}{369}{常规·球域积分}
设 $\Omega=\{x^2+y^2+z^2\leq1,x,y,z\geq0\}$。用球面坐标计算 $\iiint_\Omega xyz\,\dd V$。
\end{exercise}
\begin{exercise}{370}{基础·球锥交域}
设 $\Omega=\{x^2+y^2+z^2\leq1,z\geq\sqrt{x^2+y^2}\}$。写出球面坐标范围并求体积，极角从正 $z$ 轴起算。
\end{exercise}
\begin{exercise}{371}{基础·半球对称性}
设 $\Omega=\{x^2+y^2+z^2\leq1,z\geq0\}$。利用对称性计算 $\iiint_\Omega(x+y)\,\dd V$，再计算 $\iiint_\Omega z\,\dd V$。
\end{exercise}
\begin{exercise}{372}{基础·球域二次矩}
设 $\Omega$ 为单位球体。利用 $x,y,z$ 的对称性计算 $\iiint_\Omega(x^2+2y^2+3z^2)\,\dd V$。
\end{exercise}
\begin{exercise}{373}{基础·球体转动惯量}
半径为 $R>0$、密度为 $1$ 的实心球以原点为球心。求其关于 $z$ 轴的转动惯量。
\end{exercise}
\begin{exercise}{374}{常规·球锥域一阶矩}
设 $\Omega=\{x^2+y^2+z^2\leq1,z\geq\sqrt{x^2+y^2}\}$。计算 $\iiint_\Omega(x+z)\,\dd V$。
\end{exercise}
\begin{exercise}{375}{基础·径向函数积分}
分别计算单位球体上的 $\iiint(x^2+y^2+z^2)\,\dd V$ 与单位圆盘上的 $\iint(x^2+y^2)\,\dd\sigma$，求两者的比值。
\end{exercise}
\begin{exercise}{376}{基础·圆柱体积分}
设 $\Omega=\{x^2+y^2\leq1,0\leq z\leq2\}$。用柱面坐标计算 $\iiint_\Omega(x^2+y^2)\,\dd V$。
\end{exercise}
\begin{exercise}{377}{基础·球冠体积变式}
设 $\Omega=\{x^2+y^2+z^2\leq1,z\geq1/2\}$。用水平截面法求体积。
\end{exercise}
\begin{exercise}{378}{常规·抛物体质心小综合}
设 $\Omega=\{x^2+y^2\leq z\leq1\}$，体密度为 $\rho=1+z$。求质量及质心坐标。
\end{exercise}
\subsection{重积分的应用}
\begin{exercise}{379}{基础·平面三角形面积}
用曲面面积积分求平面 $x+y+z=1$ 在第一卦限部分的面积。
\end{exercise}
\begin{exercise}{380}{基础·圆柱侧面积}
用参数曲面的面积公式求圆柱面 $x^2+y^2=4,0\leq z\leq3$ 的侧面积。
\end{exercise}
\begin{exercise}{381}{基础·球带面积}
求单位球面上 $0\leq z\leq1/2$ 部分的面积。可用球面参数或球带公式。
\end{exercise}
\begin{exercise}{382}{常规·薄片质量与质心}
设 $D=\{0\leq x\leq1,0\leq y\leq x\}$，面密度为 $\rho=x$。求质量及质心。
\end{exercise}
\begin{exercise}{383}{基础·对称密度的质心}
设 $\Omega$ 为单位球体，体密度为 $\rho=1+x^2+y^2+z^2$。求质量，并利用对称性确定质心。
\end{exercise}
\begin{exercise}{384}{常规·抛物体形心}
设均匀物体为 $\Omega=\{x^2+y^2\leq z\leq1\}$。求其体积与形心坐标。
\end{exercise}
\begin{exercise}{385}{基础·圆锥形心}
设均匀圆锥体为 $\Omega=\{0\leq z\leq1,x^2+y^2\leq(1-z)^2\}$。利用水平截面求形心高度。
\end{exercise}
\begin{exercise}{386}{基础·椭圆薄片惯量}
设 $a,b>0$，均匀薄片 $D=\{x^2/a^2+y^2/b^2\leq1\}$ 的面密度为 $1$。求关于 $x$ 轴与 $y$ 轴的转动惯量。
\end{exercise}
\begin{exercise}{387}{基础·矩形薄片惯量}
设薄片为 $D=[-1,1]\times[0,1]$，面密度为 $1$。求关于直线 $y=1/2$ 的转动惯量。
\end{exercise}
\begin{exercise}{388}{常规·旋转体轴惯量}
设 $\Omega=\{x^2+y^2\leq z\leq1\}$，体密度为 $1$。求关于 $z$ 轴的转动惯量及惯量与质量之比。
\end{exercise}
\begin{exercise}{389}{常规·薄圆盘轴向引力}
半径为 $1$ 的均匀薄圆盘位于 $z=0$，面密度为 $1$，引力常数为 $G$。求它对 $(0,0,a)$（$a>0$）处单位质量质点的引力向量。提示：先用对称性确定方向，再用极坐标求轴向分量。
\end{exercise}
\begin{exercise}{390}{基础·上半球面积变式}
用球面参数 $x=\sin\varphi\cos\theta,y=\sin\varphi\sin\theta,z=\cos\varphi$ 求单位上半球面积，写明 $\varphi,\theta$ 的范围。
\end{exercise}
\begin{exercise}{391}{常规·圆盘非均匀密度小综合}
设 $D=\{x^2+y^2\leq1\}$，面密度为 $\rho=1+x/2$。求质量与质心坐标，并计算关于 $y$ 轴的转动惯量。
\end{exercise}
